\section{Stereographic Encoding}
\label{sec:encoding}

A single-qubit pure state
$\ket{\psi} = \psi_0\ket{0} + \psi_1\ket{1}$ lives in the
two-dimensional Hilbert space $\C^2$, subject to the normalization
$|\psi_0|^2 + |\psi_1|^2 = 1$.  Since the global phase is
unobservable, the physical state space is the complex projective
line $\C P^1$, which is topologically a sphere.  This sphere is
made concrete by the Bloch representation: the expectation values
of the three Pauli matrices define a unit vector
\begin{equation}
	\vec{n} = \big(\langle\PauliX\rangle,\;
	\langle\PauliY\rangle,\;
	\langle\PauliZ\rangle\big) \in \Sphere,
	\label{eq:bloch-vector}
\end{equation}
where $\langle\sigma_j\rangle = \bra{\psi}\sigma_j\ket{\psi}$,
with the north pole $(0,0,1)$ corresponding to $\ket{0}$ and the
south pole to $\ket{1}$.

Encoding classical data into this state space is the first step in
any quantum algorithm that processes classical inputs.  The most
common strategy, amplitude encoding, writes a real number $x$
directly into a state amplitude:
$\ket{\psi} = x\ket{0} + \sqrt{1-x^2}\,\ket{1}$, which
immediately forces $|x| \leq 1$.  Values outside $[-1,1]$ must be
rescaled by a normalization factor that propagates through the
subsequent computation---precisely the mechanism behind the
boundedness constraint of QSP discussed in
Section~\ref{sec:preliminaries}.

Stereographic projection~\cite{DubrovinFomenkoNovikov1992}
provides a fundamentally different encoding that avoids this
restriction.  Projecting from the north pole of the unit sphere
$\Sphere \subset \R^3$ onto the equatorial plane defines a
conformal bijection
\begin{equation}
	(u, v, w) \;\longmapsto\; z = \frac{u + iv}{1 - w},
	\label{eq:stereo-forward}
\end{equation}
from $\Sphere \setminus \{(0,0,1)\}$ to $\C$, with the
convention that the north pole maps to $z = \infty$ and the
south pole $(0,0,-1)$ to $z = 0$.  The inverse map sends a
complex number back to the sphere:
\begin{equation}
	z \;\longmapsto\; (u, v, w) = \left(\frac{2\re(z)}{|z|^2 + 1},\;
	\frac{2\im(z)}{|z|^2 + 1},\;
	\frac{|z|^2 - 1}{|z|^2 + 1}\right),
	\label{eq:stereo-inverse}
\end{equation}
where $(u,v,w) \in \Sphere$.  The crucial observation is that the Bloch
sphere~\eqref{eq:bloch-vector} and the stereographic
sphere~\eqref{eq:stereo-inverse} are the \emph{same} sphere: if
we identify $z$ with the amplitude ratio~$\psi_0/\psi_1$, the
Bloch vector of~$\ket{\psi}$ is exactly the stereographic image
of~$z$.  Every complex number---no matter how large---therefore
has a unique image on the Bloch sphere, with the radial coordinate
$r = |z| \in [0,\infty)$ compressed into the bounded altitude
$w = (r^2 - 1)/(r^2 + 1) \in [-1,1)$.  This motivates the
following definition.

\begin{definition}[Stereographic Encoding]
	\label{def:stereo-encoding}
	For $z \in \C$, the stereographically encoded quantum state is
	\begin{equation}
		\ket{z} = \frac{1}{\sqrt{|z|^2 + 1}}\big(z\ket{0} + \ket{1}\big).
		\label{eq:encoding}
	\end{equation}
	In polar form $z = re^{i\theta}$ with $r \geq 0$, this reads
	\begin{equation}
		\ket{z} = \frac{1}{\sqrt{r^2+1}}\big(re^{i\theta}\ket{0} + \ket{1}\big).
		\label{eq:encoding-polar}
	\end{equation}
\end{definition}

The state~\eqref{eq:encoding} is normalized for all $z \in \C$,
and the encoding maps $\C \cup \{\infty\}$ bijectively to the
Bloch sphere: $z = 0$ maps to $\ket{1}$ (south pole),
$z \to \infty$ to $\ket{0}$ (north pole), and $|z| = 1$ traces
the equator.  Computing the Bloch vector of $\ket{z}$ gives
\begin{equation}
	\left(\langle\PauliX\rangle,\,
	\langle\PauliY\rangle,\,
	\langle\PauliZ\rangle\right) =
	\left(\frac{2\re(z)}{|z|^2+1},\;
	\frac{2\im(z)}{|z|^2+1},\;
	\frac{|z|^2-1}{|z|^2+1}\right),
	\label{eq:bloch-components}
\end{equation}
confirming the identification: the Bloch vector of $\ket{z}$ is
the stereographic image of~$z$.

An encoding is only useful if the encoded data can be read back
out.  The decoding formula follows directly
from~\eqref{eq:bloch-components}: since
$\langle\PauliX\rangle + i\langle\PauliY\rangle = 2z/(|z|^2+1)$
and $1 - \langle\PauliZ\rangle = 2/(|z|^2+1)$, their ratio
recovers~$z$:
\begin{equation}
	z = \frac{\langle\PauliX\rangle + i\langle\PauliY\rangle}
	{1 - \langle\PauliZ\rangle}.
	\label{eq:decoding}
\end{equation}
This is precisely the inverse stereographic projection---the same
map that sends Bloch coordinates back to the complex plane.
Equivalently, for a state
$\alpha\ket{0} + \beta\ket{1}$ in stereographic encoding, the
decoded value is the amplitude ratio $\alpha/\beta$.  This
observation is central to our main result: QSP produces a state
$P(\rtilde)\ket{0} + Q(\rtilde)\ket{1}$, and decoding returns
the potentially unbounded ratio~$P/Q$.

The encoded state~\eqref{eq:encoding} is prepared from
$\ket{0}$ by the unitary
\begin{equation}
	U_z = \frac{1}{\sqrt{1+r^2}}\!\begin{pmatrix}
		re^{i\theta} & 1 \\ 1 & -re^{-i\theta}
	\end{pmatrix},
	\label{eq:encoding-unitary}
\end{equation}
satisfying $U_z\ket{0} = \ket{z}$ and $U_z^\dagger U_z = I$.
In terms of $z$ directly, this reads
\begin{equation}
	U_z = \frac{1}{\sqrt{1+|z|^2}}
	\begin{pmatrix} z & 1 \\ 1 & -\bar{z}\end{pmatrix},
	\label{eq:encoding-unitary-compact}
\end{equation}
whose columns are $\ket{z}$ and its orthogonal complement.  Note
that $\det(U_z) = -1$, so $U_z \in \U(2)$ but not $\SU(2)$.

With the encoding, decoding, and state preparation in hand, the
natural question is what happens when quantum gates act on an
encoded state.  Since $\ket{z}$ is parametrized by a complex
number, any unitary $U$ sends $\ket{z}$ to a new encoded state
$\ket{f(z)}$---and the resulting function $f$ turns out to be a
M\"obius transformation.

\subsection{Quantum Gates as M\"obius Transformations}
\label{subsec:mobius}

Since decoding reads off the amplitude ratio $\alpha/\beta$,
the action of any single-qubit gate on an encoded state has a
clean algebraic description.  A unitary $U$ with matrix entries
$(a,b;\,c,d)$ sends $\ket{z} \propto z\ket{0} + \ket{1}$ to
a state with amplitudes $\alpha = az+b$ and $\beta = cz+d$,
so the decoded output is the M\"obius transformation
$f(z) = (az+b)/(cz+d)$.  The Pauli gates give the simplest
examples:
\begin{equation}
	X: z \mapsto \frac{1}{z}, \qquad
	Y: z \mapsto -\frac{1}{z}, \qquad
	Z: z \mapsto -z.
	\label{eq:pauli-mobius}
\end{equation}
These act as inversions, negations, and reflections of the
encoded value---familiar M\"obius transformations of the Riemann
sphere.  More generally, the gate
$U_3(\vartheta,\varphi,\lambda) \in \U(2)$ induces the M\"obius
map $z \mapsto (az+b)/(cz+d)$ with
\begin{equation}
	\begin{pmatrix} a & b \\ c & d \end{pmatrix} =
	\begin{pmatrix}
		e^{-i\varphi}\cos\frac{\vartheta}{2} &
		-e^{-i(\lambda+\varphi)}\sin\frac{\vartheta}{2} \\[4pt]
		e^{i\lambda}\sin\frac{\vartheta}{2}  &
		\cos\frac{\vartheta}{2}
	\end{pmatrix}.
	\label{eq:U3-mobius}
\end{equation}
The gate--M\"obius correspondence is a group homomorphism
$\U(2) \to \mathrm{PGL}(2,\C)$, the projective linear
group of M\"obius transformations.
In particular, the phase gate
$\exp[i\phi\PauliZ] = \mathrm{diag}(e^{i\phi},e^{-i\phi})$
acts as $z \mapsto e^{2i\phi}z$.

This correspondence transforms the problem of analyzing quantum
circuits on encoded states into the algebraic study of composed
M\"obius maps---that is, rational functions.  Each gate
contributes a linear fractional factor, and the composition of
the full QSP circuit yields the rational function whose numerator
and denominator are the polynomials $P$ and $Q$ that appear in
the output state.  The next section makes this precise.
